\chapter{Weiteres}

\section{Laborordnung}
Bei Arbeiten im Labor des Instituts ist an einer Unterweisung in der Laborordnung teilzunehmen. Näheres klärt der Betreuer mit Ihnen ab.

\section{Präsentation}
Für die Präsentation Ihrer Arbeit verwenden Sie bitte den Folienmaster der Professur. Die aktuelle Fassung des Folienmasters erhalten Sie auch von Ihrem Betreuer.

Die Gliederung der Präsentation sollte im Wesentlichen der Gliederung der Ausarbeitung folgen, jedoch aufgrund der begrenzten Zeit die wesentlichen Aspekte (Problemstellung, Ziel, Methodik, Ergebnisse) hervorheben. 
Die Präsentation endet mit der Zusammenfassung als Übersicht.

\section{Anzahl der Exemplare, Einbandform}
Die Arbeit ist sowohl schriftlich, fest eingebunden, als auch elektronisch auf CD zu dokumentieren. Entsprechend der jeweiligen Prüfungsordnung sind die notwendigen Exemplare beim Prüfungsamt abzugeben. Es ist für jeden Prüfer ein Exemplar vorzusehen, diese werden durch das Prüfungsamt weitergeleitet.
Die elektronische Dokumentation enthält, in einer klaren Ordnerstruktur und Dateibenennung Folgendes:
\begin{itemize}
	\item Dokumentation, Word und PDF
	\item Präsentation, Zwischen, sowie Abschlusspräsentationen
	\item Quellen, insbesondere elektronische
	\item Arbeitsdokumente, Zeitplan, Anforderungslisten, Versuchsdokumentationen o.ä.
	\item Software, lauffähige Version, Quellcode und Dokumentation
	\item Werkzeuge, sofern neue Werkzeuge eingeführt wurden
\end{itemize}
Sollten Sie dabei Datumsangaben bei der Dateibezeichnung verwenden, bietet es sich an, dies in der Form \verb|YYMMDD_Dateiname.x| zu tun. So ist eine chronologische Sortierung gewährleistet.